trace identity to $f=\theta H_{\{p\}}$:
\[
 U_p^2t(\phi)=U_p^3+d(\phi)U_p,\qquad
 (t(\phi)f)[p^2]=f[p^3]+(d(\phi)f)[p].
\]
Here and in this calculation all congruences are modulo the ideal
$2\mathcal O_{\rm coef}$ before passage to the residue field.
Write $a_p=a_p(E^{(k)})$.
We check every theta convolution term in the three coefficients.
Put
\[
\begin{gathered}
 B=B_{\{p\}}=\frac{(U_p-\eta(p)V_p)g}{2},\qquad
 b_1=B[1]=\frac{a(p)}2,\quad b_p=B[p],\\
 H=H_{\{p\}}=\frac{B-b_1\Theta_1-b_p\Theta_{\{p\}}}{2}.
\end{gathered}
\]
The inherited depth is zero, and $H$ is integral by (O.11).
Also $H[0]=0$, since the cuspidal source and the unary series have
zero constant term.
The ordinary lower bound gives $b_1\in\mathcal O_{\rm coef}$.
Under the assumed absence of a nonempty $K$-split unit seed one
has the stronger assertion
\[
                         b_1\in2\mathcal O_{\rm coef}.          \tag{G.4a}
\]
Indeed, if $p$ lies in the coefficient filter, the improved bound is
$v_2a(p)\ge2$; otherwise the actual filter gives $a(p)=0$.
No condition that $p$ itself belong to that filter is imposed.

Let $\delta_B,\delta_1,\delta_p$ be the scalars by which $d(\phi)$
acts on the full weight-two blocks
$\theta B,\theta\Theta_1,\theta\Theta_{\{p\}}$, respectively.
They are integral units.  The determinant is the cyclotomic character
times the weight-two nebentypus on every cuspidal constituent and
on the Eisenstein character pair, independently of old multiplicity.
For this positive-twist source $\psi/\eta=\chi_{-1}$ and $i=0$.
The half-integral-weight characters of $B$ and $\Theta_{\{p\}}$
are therefore both
\[
                 \psi(p/\cdot)=\eta(-p/\cdot).
\]
Multiplication by the same $\theta$ preserves this equality.
Consequently $\delta_B=\delta_p$ for the actual Frobenius lift
$\phi$, including its action on the ramified character.  The ratio
of the characters of $\theta\Theta_{\{p\}}$ and
$\theta\Theta_1$ is quadratic.  Their cyclotomic determinant factors
are identical, so $\delta_p/\delta_1\in\{1,-1\}$ and
$\delta_p-\delta_1\in2\mathcal O_{\rm coef}$.  We must not simply
equate these two scalars.  On the actual characteristic-zero
module the exact identity is
\[
 d(\phi)f=\delta_p f+
       \frac{(\delta_p-\delta_1)b_1}{2}\,\theta\Theta_1.
\]
Now $H[p]=0$, since $B[p]=b_p$, $\Theta_{\{p\}}[p]=1$ and
$\Theta_1[p]=0$.  Integrality of $H$ implies $f[p]\in2\mathcal O_{\rm coef}$.
Also $\theta\Theta_1[p]=2\sum_{0<j^2<p}\Theta_1[p-j^2]$ is even.
The displayed determinant identity therefore already proves
$(d(\phi)f)[p]\in2\mathcal O_{\rm coef}$ using only the ordinary
integrality of $b_1$.  The following direct expansion checks the
individual cross terms under the no-seed assumption in force here.

For $n_j=p-j^2>0$, one has $p\nmid n_j$ and
$\Theta_{\{p\}}[n_j]=0$.  We claim that
\[
                           B[n_j]\in2\mathcal O_{\rm coef}.    \tag{G.4b}
\]
If $n_j$ meets the old support, the coprimality sieve makes this
coefficient zero.  In particular $2\in S$, so every odd $j$ gives
an even $n_j$ and vanishes in this way.  Otherwise write
$n_j=r m^2$, with $r$ squarefree.
Since $p\nmid n_j$, the $V_p$ term is zero and
$B[n_j]=a(prm^2)/2$.  The integral good-prime shell recurrences
express $a(prm^2)$ as an integral multiple, up to unit character
factors, of $a(pr)$.  If $r>1$, the ordinary coefficient lower
bound gives $v_2a(pr)\ge\omega(pr)=1+\omega(r)\ge2$, or the
coefficient is zero off the filter.  This uses no assertion that
the primes dividing $r$ split in $K$.  If $r=1$, the same shell
recurrences express $B[m^2]$ as an integral multiple of $b_1$,
and (G.4a) applies.  This proves (G.4b).  Also every coefficient
of $b_1\Theta_1$ is even by (G.4a).

Now expand $\theta=1+2\sum_{j>0}q^{j^2}$ before discarding any
term.  Linearity of the determinant operator on the three actual
blocks gives the exact expression
\[
\begin{split}
 (d(\phi)f)[p]
 ={}&\frac{\delta_B B[p]-\delta_p b_p}{2}\\
 &+\sum_{0<j^2<p}
 \bigl(\delta_B B[p-j^2]
       -\delta_1 b_1\Theta_1[p-j^2]
       -\delta_p b_p\Theta_{\{p\}}[p-j^2]\bigr).
\end{split}                                                     \tag{G.4c}
\]
There is no constant-theta $\Theta_1[p]$ term since $p$ is not
a square, and $\Theta_{\{p\}}[p]=1$.
The first line is zero because $B[p]=b_p$ and
$\delta_B=\delta_p$.  Every summand in the second line is even
by (G.4a)--(G.4b) and the vanishing of its last term.  In particular
representations $p=j^2+m^2$, which can occur when $p\equiv1\pmod4$,
are covered by the case $r=1$ and are not omitted.  Thus
$(d(\phi)f)[p]\in2\mathcal O_{\rm coef}$.

